% ===== 4. 実験 (Experiments) =====
% 役割:
%   - HA48 アノテーションデータでの予備検証
%   - ベースライン比較 (G-Eval / BERTScore 最低 1 本)
%   - 数値の正直さ: 公式論文と docs/validation.md の同期を保つ
%
% 分量目安: 2 ページ
%
% TODO: 並行 backlog 完了後に最終数値を埋める
%   - HA48 → HA100 拡充版での再評価
%   - 第二 annotator + IAA (Cohen's κ)
%   - ベースライン (G-Eval) 比較

\section{実験}
\label{sec:experiments}

% ----- 4.1 設定 -----
\subsection{実験設定}
\label{sec:exp:setup}

% TODO: 4.1.1 データセット
% HA48: 48 件の人手アノテーション付き LLM 応答
% (拡充後は HA100 に差し替え)
\subsubsection{データセット}
% [PLACEHOLDER]
% data/human_annotation_48/annotation_48_merged.csv

% TODO: 4.1.2 比較手法 (ベースライン)
\subsubsection{比較手法}
% [PLACEHOLDER]
% - BERTScore
% - G-Eval (LLM-as-judge)
% - 提案手法 (PoR/ΔE)

% TODO: 4.1.3 評価指標
\subsubsection{評価指標}
% [PLACEHOLDER]
% - 人手評価との相関 (Spearman / Pearson)
% - verdict 分布と人手 verdict の一致率
% - Cohen's κ (annotator 間)

% ----- 4.2 結果 -----
\subsection{結果}
\label{sec:exp:results}

% TODO: 表 1 - 主要結果 (相関係数)
% docs/validation.md の current snapshot から転記

% TODO: 表 2 - verdict 分布の混同行列

% TODO: 図 1 - PoR 座標上での応答プロット

% ----- 4.3 考察 -----
\subsection{結果の解釈}
\label{sec:exp:discussion}

% TODO: L_sem 分解の効用 (どの損失項が verdict に効くか)
% TODO: 既存指標との差分 (S/C 分離の意義)
